\section*{الفصل \ref{ch:g12:age-of-the-earth} --- التأريخ الإشعاعي وعمر الأرض}

\begin{solution}{exo:g12:age-of-the-earth:1}
(أ) \omterm{def:g12:age-of-the-earth:closed-system}{نظام مغلق}؛ (ب) \omterm{def:g12:age-of-the-earth:isochron}{خط تساوي العمر}؛ (ج) \omterm{def:g12:age-of-the-earth:age}{عمر الأرض}.
\end{solution}

\begin{solution}{exo:g12:age-of-the-earth:2}
البلاجيوكليز (0.35)، فالهورنبلند (1.20)، فالصخر الكلي (3.10)، ففلسبار البوتاسيوم (6.40)،
فالبيوتيت (24.0). والبيوتيت يقع في الأعلى: فهو يحوي أكبر قدر من الروبيديوم-87 بالنسبة إلى
ما فيه من السترونتيوم-86، ولذلك صنع روبيديومه أكبر قدر من السترونتيوم-87 منذ
تكوّن الغرانيت.
\end{solution}

\begin{solution}{exo:g12:age-of-the-earth:3}
$T = \ln 2/\lambda = 0.6931/\num{1.3972e-11} \approx \num{4.96e10}$ سنة،
أي 49.6 مليار سنة: وهذا $49.6/4.54 \approx 11$ ضعفًا \omterm{def:g12:age-of-the-earth:age}{لعمر الأرض}.
\end{solution}

\begin{solution}{exo:g12:age-of-the-earth:4}
$\lambda t = \num{1.3972e-11} \times \num{3.10e8} \approx 0.004331$،
و$e^{0.004331} - 1 \approx 0.004341$. وعند 480 مليون سنة، $\lambda t
\approx 0.006707$ و$e^{\lambda t} - 1 \approx 0.006729$. وكلتا القيمتين قريبة من
$\lambda t$ نفسه، لأن $\lambda t$ صغير.
\end{solution}

\begin{solution}{exo:g12:age-of-the-earth:5}
فيليبس (1860)، فكلفن (1862)، فجولي (1899)، فهولمز (1927)، فباترسون
(1956). $4540/100 \approx 45$ ضعفًا.
\end{solution}

\begin{solution}{exo:g12:age-of-the-earth:6}
$0.7080 + 10.0 \times 0.004341 \approx 0.7514$ اليوم؛ و0.7080 حين تكوّن
الغرانيت، شأنه شأن كل معدن في الغرانيت.
\end{solution}

\begin{solution}{exo:g12:age-of-the-earth:7}
(1) $m = (0.77525 - 0.71806)/(9.40 - 0.90) = 0.05719/8.50 \approx 0.006728$؛
$t = \ln(1.006728)/\num{1.3972e-11} \approx \num{4.80e8}$ سنة، أي 480 مليون
سنة. (2) $0.71806 + 1.70 \times 0.006728 \approx 0.72950$ و$0.71806 + 4.20
\times 0.006728 \approx 0.74632$: كلتاهما على الخط. (3) نعم: فالغرانيت يقطع
الشيست، فهو أحدث منه، و310 أقل من 480 مليون سنة.
\end{solution}

\begin{solution}{exo:g12:age-of-the-earth:8}
كانت نقاطه كلها ستقع في الموضع نفسه: والنقطة الواحدة لا تحدد خطًا ولا
ميلًا. والبيوتيت يقع بعيدًا إلى يمين سائر المعادن؛ فالخطأ الصغير في
أي نقطة يغيّر حينئذ ميل الخط المار به وبالبلاجيوكليز
أقل بكثير مما يغيّره لو كانت النقطتان متقاربتين.
\end{solution}

\begin{solution}{exo:g12:age-of-the-earth:9}
(1) روبيديومه لم يتغير، فيحتفظ بقيمة $x$؛ وصار فيه
سترونتيوم-87 أقل، فتنخفض قيمة $y$: فتنتقل النقطة إلى الأسفل مباشرة. (2) الخط
المار بها وبالبلاجيوكليز أقل انحدارًا، فيعطي عمرًا أصغر من الحقيقي.
(3) لم يعد البيوتيت يقع على الخط المار بسائر المعادن: فالنقاط
غير مصطفة، وهذه علامة على أن إحدى العينات لم تكن \omterm{def:g12:age-of-the-earth:closed-system}{نظامًا مغلقًا}.
\end{solution}

\begin{solution}{exo:g12:age-of-the-earth:10}
(1) $4540/96 \approx 47$ مرة. (2) إذا غادر شيء من الملح البحر، مدفونًا
مثلًا في الرسوبيات، كان في البحر ملح أقل مما جلبته الأنهار في المجموع:
فقسمة الملح الموجود على الإمداد السنوي تعطي زمنًا أقصر من
الحقيقي.
\end{solution}

\begin{solution}{exo:g12:age-of-the-earth:11}
$m = (34.86 - 10.34)/(50.28 - 9.46) = 24.52/40.82 \approx 0.601$، بين
ميلَي 4.55 (0.590) و4.60 مليار سنة (0.610): أي نحو 4.58
مليار سنة.
\end{solution}

\begin{solution}{exo:g12:age-of-the-earth:12}
(1) البازلت، ذو النسبة الأدنى، من الوشاح؛ والغرانيت، ذو النسبة
الأعلى، من قشرة منصهرة ولو جزئيًا. (2) الروبيديوم يذهب إلى
المصهور حين ينصهر الوشاح، ولذلك تركّز في القشرة، التي
صنعت على مر العصور قدرًا أكبر من السترونتيوم-87. (3) لا: فالمعادن التي تتبلور
من صهارة واحدة تشترك في نسبتها الابتدائية، وهذان الصخران لا يشتركان فيها.
\end{solution}

\begin{solution}{exo:g12:age-of-the-earth:13}
(1) الروبيديوم-87: $2^{-4.54/49.61} \approx 0.94$، أي يبقى نحو \qty{94}{\percent}؛
واليورانيوم-238: $2^{-4.54/4.47} \approx 0.49$، أي نحو النصف. (2) $m =
e^{\lambda t} - 1 \approx \num{1.3972e-11} \times 10^6 \approx
\num{1.4e-5}$: فحتى بيوتيت غرانيت فاردن ($x = 24$) لن يرتفع
إلا بمقدار $0.0003$ فوق النسبة الابتدائية. فما تكوّن من السترونتيوم-87 قليل جدًا
حتى يصعب تمييز الميل من الصفر: فالساعة البطيئة تلائم الصخور القديمة.
\end{solution}

\begin{solution}{exo:g12:age-of-the-earth:14}
(1) $4.438 + 0.6026 \times 19.0 \approx 15.89$؛ و$15.5$ أدنى منها بمقدار $0.39$،
أي نحو \qty{2.4}{\percent}: فالرسوبيات تقع قريبًا من الخط. (2) 0.6026
تقع بين 0.590 (4.55) و0.610 (4.60): أي نحو 4.58 مليار سنة، ضمن
مدى باترسون من 4.48 إلى 4.62 مليار سنة.
\end{solution}

\begin{solution}{exo:g12:age-of-the-earth:15}
(1) لا يفقد المخزون حينئذ إلا $46.7 - 24.4 = \qty{22.3}{TW}$ صافيًا، فيدوم
المخزون نفسه $46.7/22.3 \approx 2.1$ مرة أطول. (2) $96 \times 2.1 \approx
200$ مليون سنة، وهذا لا يزال أقصر من الحقيقة بأكثر من عشرين مرة. (3) \omterm{def:g12:internal-heat:convection}{الحمل
الحراري في الوشاح}، الذي يحمل الحرارة صاعدًا من الأعماق ويُبقي التدرج
قرب السطح مرتفعًا مدة أطول بكثير مما يفعله التوصيل وحده.
\end{solution}

\begin{solution}{pb:g12:age-of-the-earth:1}
\textbf{1.} البيوتيت ($x = 24.0$).

\textbf{2.} تبلورت من صهارة واحدة، والمعادن التي تتبلور
من مصهور واحد تأخذ سترونتيومه بنسبة نظائره نفسها.

\textbf{3.} السترونتيوم-86 مستقر ولا يصنعه أي تفكك: ففي \omterm{def:g12:age-of-the-earth:closed-system}{النظام المغلق}
لا يتغير عدده، فهو مقياس مرجعي ثابت.

\textbf{4.} السترونتيوم-87، الذي يصنعه تفكك الروبيديوم-87؛ أما الروبيديوم-87 فقد
تناقص، وبقي السترونتيوم-86 على حاله.

\textbf{5.} خطًا أفقيًا عند $y = 0.7080$: النسبة نفسها، وقيم $x$ مختلفة.

\textbf{6.} $(0.81218 - 0.70952)/(24.0 - 0.35) = 0.10266/23.65 \approx
0.004341$.

\textbf{7.} $(0.72146 - 0.70952)/(3.10 - 0.35) = 0.01194/2.75 \approx
0.004342$.

\textbf{8.} يتفقان في حدود واحد من 4000. والأول أوثق:
فنقطتاه متباعدتان، فيقل أثر تقريب النسب.

\textbf{9.} $0.70952 - 0.35 \times 0.004341 \approx 0.7080$.

\textbf{10.} $0.7080 + 6.40 \times 0.004341 \approx 0.73578$: وهي
القيمة المقيسة تمامًا.

\textbf{11.} $e^{\lambda t} = 1 + m = 1.004341$.

\textbf{12.} $\lambda t = \ln 1.004341 \approx 0.004331$.

\textbf{13.} $t = 0.004331/\num{1.3972e-11} \approx \num{3.10e8}$ سنة.

\textbf{14.} نعم: فالغرانيت يقطع الشيست، فهو بحسب مبدأ
التقاطع أحدث منه، و310 أقل من 480 مليون
سنة.

\textbf{15.} تقع 310 ملايين سنة بين 358.9 و298.9 مليون سنة:
أي في الكربوني.

\textbf{16.} نسبة $^{87}\mathrm{Sr}/^{86}\mathrm{Sr}$ في الصهارة حين
تبلور الغرانيت.

\textbf{17.} الغرانيت: $0.7080 > 0.7035$.

\textbf{18.} الروبيديوم يذهب إلى المصهور حين ينصهر الوشاح، ولذلك
تركّز في القشرة؛ ولأن الروبيديوم فيها أكثر، صنعت القشرة قدرًا أكبر من
السترونتيوم-87.

\textbf{19.} أن الصهارة جاءت، جزئيًا على الأقل، من \omterm{def:g12:continental-crust:continental-crust}{قشرة قارية} منصهرة
لا من الوشاح وحده.

\textbf{20.} عمر غرانيت فاردن \textbf{310 ملايين سنة}.
\end{solution}
